\documentclass[
    preprint,
%    twocolumn,
    linenumbers,
    amsmath,amssymb,amsfonts,
    aps, physrev,
    superscriptaddress,
%    floatfix,
%    draft,
    10pt,
    bibnotes
    ]{revtex4-2}

\usepackage{bm}
\usepackage[T1]{fontenc}
\usepackage{inputenc}
\setcounter{secnumdepth}{3}
\usepackage{upgreek,babel,calc,mathrsfs,textgreek,bbm}
\usepackage{amsmath,amssymb,mathtools,graphicx,xcolor}


\newcommand{\erf}{\text{erf}}
\newcommand{\mt}[1]{{\color{green}{[MT: #1]}}}
\newcommand{\ld}[1]{{\color{red}{[LD: #1]}}}

\bibliographystyle{apsrev4-2}

\usepackage[unicode=true,pdfusetitle,
bookmarks=true,bookmarksnumbered=false,bookmarksopen=false,
breaklinks=false,pdfborder={0 0 1},backref=false,colorlinks=true]
{hyperref}


\begin{document}

\title{Random generalized Lotka-Volterra model on sparse graphs:\\
Gaussian expansion of dynamic cavity method}

\author{Mattia Tarabolo}
\email{mattia.tarabolo@polito.it}
\affiliation{Institute of Condensed Matter Physics and Complex Systems, Department of Applied Science and Technology,
Politecnico di Torino, C.so Duca degli Abruzzi 24, I-10129 Torino, Italy}

\author{Luca Dall'Asta}
\affiliation{Institute of Condensed Matter Physics and Complex Systems, Department of Applied Science and Technology,
Politecnico di Torino, C.so Duca degli Abruzzi 24, I-10129 Torino, Italy}
\affiliation{Italian Institute for Genomic Medicine (IIGM) and Candiolo Cancer Institute IRCCS, str. prov. 142, km 3.95, I-10060 Candiolo (TO), Italy}
\affiliation{INFN, Turin Via Pietro Giuria 1, I-10125 Turin, Italy}

\author{Roberto Mulet}
\affiliation{Group of Complex Systems and Statistical Physics, Department of Theoretical Physics, Physics Faculty, University of Havana, Cuba}

\author{Alberto R. Batista-Tomás}
\affiliation{Group of Complex Systems and Statistical Physics, Department of Theoretical Physics, Physics Faculty, University of Havana, Cuba}

\date{\today}

\begin{abstract}
An article usually includes an abstract, a concise summary of the
work covered at length in the main body of the article. 
\end{abstract}

\maketitle

\section{\label{sec:intro}Introduction}

Intro on the subject.

\section{Method}

Explain here the small-coupling method and population dynamics.

\section{Results}

Present here reu


\subsection{Random generalized Lotka-Volterra model on sparse graphs}

We consider a generalized Lotka-Volterra (gLV) model describing the dynamics
of an interacting community of $N$ species, labeled by $i=1,\dots,N$.
The species interacts on an undirected random graph $G$ with symmetric
adjacency matrix $A=\{a_{ij}\}$, meaning that species $i$ and $j$
are interacting if $a_{ij}=a_{ji}=1$, otherwise the matrix elements
are zero. The time-dependent number density of individuals of species
$i$ is denoted by $x_{i}(t)$, and evolves in time according to

\begin{equation}
\frac{d}{dt}x_{i}(t)=x_{i}(t)\left[1-x_{i}(t)+\sum_{j\neq i}a_{ij}J_{ij}x_{j}(t)\right]+\eta_{i}(t) + \lambda,\label{eq:Lotka-Volterra_eq}
\end{equation}

where we included, for the sake of generality, a Gaussian noise $\eta_{i}(t)$ with vanishing mean and variance $\langle\eta_{i}(t)\eta_{i'}(t')\rangle=2Tg(x_i(t))\delta_{ii'}\delta(t-t')$, and a constant immigration rate $\lambda$. We will consider in what follows three choices for the function $g(x)$: (i) $g(x)=0$ for deterministic gLV model, (ii) $g(x)=x$ for demographic noise, and (iii) $g(x)=x^2$ for environmental noise.

The couplings $J_{ij}$ represent the (per capita) effect of species
on one another. A negative coefficient $J_{ij}$ indicates a competitive
effect of species $j$ on species $i$.

We consider here Gaussian $J_{ij}$ with average $\overline{J_{ij}}=\overline{J_{ji}}=m/K$,
variance $\text{Var}(J_{ij})=\text{Var}(J_{ji})=\sigma^{2}/K$ and
covariance $\text{Corr}(J_{ij},J_{ji})=\gamma$. The overbar $\overline{\dots}$
represents the average over the random couplings ensemble. The scaling
of the moments with the average degree of the graph $K$ is needed
to produce a well-defined fully connected limit $K\to\infty$. The
parameter $-1\leq\gamma\leq1$, being related to the correlation between
the interaction coefficients $J_{ij}$ and $J_{ji}$, measures the
fraction of predator-prey pairs. Such a predator-prey pair consists
of two species $i$ and $j$ for which $J_{ij}$ and $J_{ji}$ have
opposite signs. If for example $J_{ij}$ is positive and $J_{ji}$
is negative, the presence of species $i$ has a detrimental effect
on species $j$, while the presence of species $j$ is beneficial
for species $i$. For any pair $i<j$ we can write

\begin{equation}
J_{ij}=\frac{m}{K}+\frac{\sigma}{\sqrt{K}}z_{ij},\qquad J_{ji}=\frac{m}{K}+\frac{\sigma}{\sqrt{K}}z_{ji},
\end{equation}

where $z_{ij}$ and $z_{ji}$ are Gaussian random variables with $\overline{z_{ij}}=0$,
$\overline{z_{ij}^{2}}=1$, $\overline{z_{ij}z_{ji}}=\gamma$.

For $\gamma=1$ we have $z_{ij}=z_{ji}$ with probability one, for
$\gamma=0$ $z_{ij}$ and $z_{ji}$ are uncorrelated, and for $\gamma=-1$
we have $z_{ij}=z_{ji}$ with probability one. The latter, $\gamma=-1$,
corresponds to having all the interactions of the predator-prey type.
Increasing the value of $\gamma$ this fraction decrease, until the
opposite case , $\gamma=1$ of no predator-prey type interactions.
The other interactions can be of a mutualistic type ($J_{ij}$ and
$J_{ji}$ both positive) or strictly competitive ($J_{ij}$ and $J_{ji}$
both negative).

\subsection{Small coupling expansion of dynamic cavity equations}

We use a small coupling expansion of the dynamic cavity equations in a way similar to \cite{taraboloGaussianApproximationDynamic2025}
to derive an effective Langevin equation for non-interactive species. This approach
begins by constructing the Martin-Siggia-Rose-Janssen-De Dominicis
(MSRJD) dynamic partition function through path integrals of the trajectories
\cite{martinStatisticalDynamicsClassical1973,janssenLagrangeanClassicalField1976,dedominicisDynamicsSubstituteReplicas1978}


\begin{equation}
Z=\left\langle \int\mathcal{D}\vec{x}\prod_{i}p_{i}\left(x_{i}\left(0\right)\right)\delta\left(\dot{x_{i}}(t)-x_{i}(t)\left[1-x_{i}(t)+\sum_{j\neq i}a_{ij}J_{ij}x_{j}(t)+h_{i}(t)\right]-\eta_{i}(t)-\lambda\right)\right\rangle _{\vec{\eta}},
\end{equation}

where $\int\mathcal{D}\vec{x}=\prod_{i}\int\mathcal{D}x_{i}$ is a functional integral over the trajectories of the system, and $h_{i}(t)$ is an external field, needed to compute response functions. Following the steps of \cite{taraboloGaussianApproximationDynamic2025} we put forward a dynamic cavity Ansatz and expand the cavity marginals for small couplings. This results in a set of effective de-coupled Langevin equations on the cavity graph. For any edge $(i,j)$ on the graph $G$ the effective cavity population size $x_{i\setminus j}(t)$ evolves under the following stochastic process
\begin{align}
  \frac{d}{dt}x_{i\setminus j}(t) &= x_{i\setminus j}(t)\Biggl[1-x_{i\setminus j}(t)+\sum_{k\in\partial i\setminus j}J_{ik}\mu_{k\setminus i}(t) + \sum_{k\in\partial i\setminus j}\int_{0}^{t}dt'J_{ik}J_{ki}R_{k\setminus i}(t,t')x_{i\setminus j}(t') +h_{i\setminus j}(t)+\zeta_{i\setminus j}(t)\Biggr]\nonumber \\
  &\qquad+\sqrt{2Tg(x_{i\setminus j}(t))}\xi_{i\setminus j}(t)+\lambda\label{eq:eff_proc_cav}
\end{align}

where $h_{i\setminus j}(t)$ is an external field, needed to compute response functions, and $\zeta_{i\setminus j}$ and $\xi_{i\setminus j}$ are Gaussian noises with vanishing mean and correlations

\begin{equation}
  \left\langle \zeta_{i\setminus j}(t)\zeta_{i\setminus j}(t')\right\rangle _{i\setminus j}=\sum_{k\in\partial i\setminus j}J_{ik}^{2}C_{k\setminus i}(t,t'),\qquad\left\langle \xi_{i\setminus j}(t)\xi_{i\setminus j}(t')\right\rangle _{i\setminus j}=\delta(t-t').
\end{equation}

The terms $\mu_{k\setminus i}(t)$, $R_{k\setminus i}(t,t')$, and $C_{k\setminus i}(t,t')$ are respectively the cavity mean abundance, the cavity response function, and the cavity (connected) two-point correlation function. They are obtained self-consistently from the effective dynamics of the neighbors

\begin{equation}
\mu_{k\setminus i}(t)=\left\langle x_{k\setminus i}(t)\right\rangle _{k\setminus i},\quad R_{k\setminus i}(t,t')=\frac{\delta}{\delta h_{k\setminus i}(t')}\left\langle x_{k\setminus i}(t)\right\rangle _{k\setminus i},\quad C(t,t')=\left\langle x_{k\setminus i}(t)x_{k\setminus i}(t')\right\rangle _{k\setminus i}-\mu_{k\setminus i}(t)\mu_{k\setminus i}(t'),
\end{equation}

where $\left\langle \dots\right\rangle _{k\setminus i}$ describes the average over the distribution of the effective process Eq. \eqref{eq:eff_proc_cav} for the edge $(k,i)$. The effective process is therefore de-coupled, and the interactions are captured by the self-consistent external fields $\mu_{k\setminus i}(t)$, $R_{k\setminus i}(t,t')$, and $C_{k\setminus i}(t,t')$. A detailed derivation can be found in Appendix \ref{sec:effective_cav_proc}.

Once the cavity averages have been determined through self-consistency, one can obtain the effective de-coupled Langevin equations for the species population

\begin{equation}
  \frac{d}{dt}x_{i}(t) =x_{i}(t)\Biggl[1-x_{i}(t)+\sum_{k\in\partial i}J_{ik}\mu_{k\setminus i}(t) + \sum_{k\in\partial i}\int_{0}^{t}dt'J_{ik}J_{ki}R_{k\setminus i}(t,t')x_{i}(t')
+h_{i}(t)+\zeta_{i}(t)\Biggr]+\sqrt{2Tg(x_{i}(t))}\xi_{i}(t)+\lambda,\label{eq:eff_proc_marg}
\end{equation}

with Gaussian noises with vanishing mean and correlations

\begin{equation}
  \left\langle \zeta_{i}(t)\zeta_{i}(t')\right\rangle _{i}=\sum_{k\in\partial i}J_{ik}^{2}C_{k\setminus i}(t,t'),\qquad\left\langle \xi_{i}(t)\xi_{i}(t')\right\rangle _{i}=\delta(t-t').
\end{equation}

The average $\left\langle \dots\right\rangle _{i}$ is performed over the effective de-coupled process Eq. \eqref{eq:eff_proc_marg}.

\subsection{Numerical integration}
We propose a numerical integration scheme, mostly based on what has been proposed in \cite{royNumericalImplementationDynamical2019}. The algorithm iteratively updates the self-consistent external fields until convergence.

We start with initial guesses for the cavity mean abundances $\{\mu_{i\setminus j}\}_{(i,j)\in E}$, the cavity correlations $\{C_{i\setminus j}\}_{(i,j)\in E}$, and the cavity responses $\{R_{i\setminus j}\}_{(i,j)\in E}$ for each edge of the graph. At iteration $n$, for each edge $(i,j)$ the following steps are followed:
\begin{enumerate}
    \item We sample $N_{\rm traj}$ Gaussian paths $(\zeta_{i\setminus t}(t), \xi_{i\setminus j}(t))$, where $\zeta_{i\setminus t}(t)$ and $\xi_{i\setminus t}(t)$ have, respectively, covariance matrices $C^{\zeta}_{i\setminus j}(t,t')=\sum_{k\in\partial i \setminus j}J_{ik}^2C_{k\setminus i}(t,t')$ and $C^{\xi}_{i\setminus j}(t,t')=\delta(t,t')$.
    \item For each path, we numerically integrate the effective cavity equation with an initial condition sampled according to $p_0(x)$:
    \begin{align*}
          \frac{d}{dt}x_{i\setminus j}(t) & =x_{i\setminus j}(t)\Biggl[1-x_{i\setminus j}(t)+\sum_{k\in\partial i\setminus j}J_{ik}\mu_{k\setminus i}(t) + \sum_{k\in\partial i\setminus j}\int_{0}^{t}dt'J_{ik}J_{ki}R_{k\setminus i}(t,t')x_{i\setminus j}(t')+\zeta_{i\setminus j}(t)\Biggr]\\
          &\qquad+\sqrt{2Tg(x_{i\setminus j}(t))}\xi_{i\setminus j}(t)+\lambda,\\
          x_{i\setminus j}(0) & \sim p_0(x_{i\setminus j}(0)).
    \end{align*}
    \item We compute the new values of the cavity self-consistent fields as averages over the paths:
    \begin{align*}
        \mu^{\rm new}_{i\setminus j}(t) & = \mathbb{E}_{\rm paths}\left[x_{i\setminus j}(t) \right],\\
        C^{\rm new}_{i\setminus j}(t,t') & = \mathbb{E}_{\rm paths}\left[x_{i\setminus j}(t)x_{i\setminus j}(t') \right] - \mu^{\rm new}_{i\setminus j}(t)\mu^{\rm new}_{i\setminus j}(t'),\\
        R^{\rm new}_{i\setminus j}(t,t') & = \mathbb{E}_{\rm paths}\left[R^{\rm sp}_{i\setminus j}(t,t') \right].
    \end{align*}
    \item We update softly the cavity self-consistent fields: $X_{i\setminus j}^{\rm updated}=(1-d)X_{i\setminus j}^{\rm new}+dX_{i\setminus j}$, where $X$ is $\mu$, $C$ and $R$. $d$ is the damping factor.
\end{enumerate}

The terms $R^{\rm sp}_{i\setminus j}(t,t')$ are single-path response functions, and can be computed through numerical integration of the equation obtained by directly applying $\delta / \delta h_{i\setminus j}(t')$ to Eq.~\eqref{eq:eff_proc_cav}. More details can be found in Appendix~\ref{app:response}.

The algorithm proceeds in iterations, where each iteration consists of updating all edges $
(i,j)$ of the graph. These iterations repeat until convergence is achieved, defined as the change in the cavity self-consistent fields between successive iterations dropping below a threshold $\varepsilon$. For computational efficiency, we evaluate convergence based on the stabilization of correlation functions only.
\begin{equation}
    \max_{(i,j)}\left\lVert C^{\rm updated}_{i\setminus j} - C_{i\setminus j}^{\rm new} \right\rVert < \epsilon,
\end{equation}
where $\lVert \cdots \rVert$ is some norm to be defined. When convergence is reached, we can compute the local effective process by using the self-consistent cavity fields in Eq. \eqref{eq:eff_proc_marg}. In practice, for each node $i$ of the graph we sample $N_{\rm traj}$ Gaussian paths $(\zeta_{i}(t), \xi_{i}(t))$, where $\zeta_{i}(t)$ and $\xi_{i}(t)$ have, respectively, covariance matrices $C^{\zeta}_{i}(t,t')=\sum_{k\in\partial i}J_{ik}^2C_{k\setminus i}(t,t')$ and $C^{\xi}_{i}(t,t')=\delta(t,t')$. We then numerically integrate the effective marginal equation with an initial condition sampled according to $p_0(x)$:
\begin{align*}
      \frac{d}{dt}x_{i}(t) & =x_{i}(t)\Biggl[1-x_{i}(t)+\sum_{k\in\partial i}J_{ik}\mu_{k\setminus i}(t) + \sum_{k\in\partial i}\int_{0}^{t}dt'J_{ik}J_{ki}R_{k\setminus i}(t,t')x_{i}(t')+\zeta_{i}(t)\Biggr]+\sqrt{2Tg(x_{i}(t))}\xi_{i}(t)+\lambda,\\
      x_{i}(0) & \sim p_0(x_{i}(0)). 
\end{align*}
The full average quantities are then obtained by averaging over the sampled paths:
\begin{align*}
    \mu_{i}(t) & = \mathbb{E}_{\rm paths}\left[x_{i}(t) \right],\\
    C_{i}(t,t') & = \mathbb{E}_{\rm paths}\left[x_{i}(t)x_{i}(t') \right] - \mu_{i}(t)\mu_{i}(t'),\\
    R_{i}(t,t') & = \mathbb{E}_{\rm paths}\left[R^{\rm sp}_{i}(t,t') \right].
\end{align*}


More details on the numerical integration can be found in Appendix~\ref{app:numerical_integration}.

    
\subsection{Gaussian closure approximation}
While a direct numerical integration of the effective process is possible, we can also use a Gaussian closure approximation to obtain analytical results. This approach is based on the assumption that the distribution of the population sizes $x_i(t)$ is Gaussian, which allows us to compute the moments of the process analytically. The Gaussian closure approximation leads to a set of self-consistent equations for the first and second moments of the population sizes, which can be solved iteratively.

In particular, by taking averages of the effective cavity process Eq. \eqref{eq:eff_proc_cav} and approximating three-times averages as sums of products of one- and two-times averages, we obtain the following self-consistent equations for the first and second moments of the population sizes:
\begin{align}
  \frac{d}{dt}\mu_{i\setminus j}(t) & = \langle \dot{x}_{i\setminus j}(t) \rangle_{i\setminus j} = \mu_{i\setminus j}(t)-C_{i\setminus j}(t,t)-\mu_{i\setminus j}^2(t)+\sum_{k\in\partial i \setminus j}J_{ik}\mu_{k\setminus i}(t)\mu_{i\setminus j}(t)  \nonumber \\
  & \quad +\sum_{k\in \partial i \setminus j}J_{ik}J_{ki}\int_{0}^{t}ds\, R_{k\setminus i}(t, s)\left(C_{i\setminus j}(t, s)+\mu_{i\setminus j}(t)\mu_{i\setminus j}(s)\right) + \lambda,\\
  \frac{\partial}{\partial t}C_{i\setminus j}(t,t') & = \langle \dot{x}_{i\setminus j}(t) x_{i\setminus j}(t') \rangle_{i\setminus j} = \left(1 - 2 \mu_{i\setminus j}(t) \right) C_{i\setminus j}(t,t') + \sum_{k\in\partial i \setminus j}J_{ik}\mu_{k\setminus i}(t)C_{i\setminus j}(t,t') \nonumber\\
  & \quad + \sum_{k\in \partial i \setminus j}J_{ik}J_{ki}\int_0^t ds\, R_{k\setminus i}(t,s)\left(\mu_{i\setminus j}(t)C_{i\setminus j}(s,t') + C_{i\setminus j}(t,t')\mu_{i\setminus j}(s)\right) + 2T \mathcal{L}_{i\setminus j}(t)R_{i\setminus j}(t',t) \nonumber \\
  & \quad  + \sum_{k\in \partial i \setminus j} J_{ik}^2 \left(\int_0^{t'} ds \, C_{k\setminus i}(t,s) \mu_{i\setminus j}(t) R_{i\setminus j}(t',s) + \int_0^{t} ds \, C_{k\setminus i}(t,s) R_{i\setminus j}(t,s) \mu_{i\setminus j}(t') \right),\\
  \frac{\partial}{\partial t}R_{i\setminus j}(t,t') & =  \langle \dot{x}_{i\setminus j}(t) {\rm i} \hat{x}_{i\setminus j}(t') \rangle_{i\setminus j} = \left(1 - 2 \mu_{i\setminus j}(t) \right) R_{i\setminus j}(t,t') + \sum_{k\in\partial i \setminus j}J_{ik}\mu_{k\setminus i}(t)R_{i\setminus j}(t,t') + \delta(t-t') \nonumber\\
  & \quad + \sum_{k\in \partial i \setminus j} J_{ik}J_{ki} \left( \int_{0}^{t} ds \, R_{k\setminus i}(t,s)R_{i\setminus j}(t,t')\mu_{i\setminus j}(s)+ \int_{t'}^{t} ds \, \mu_{i\setminus j}(t)R_{k\setminus i}(t,s) R_{i\setminus j}(s,t') \right),
\end{align}
where we have introduced the function
\begin{equation}
  \mathcal{L}_{i\setminus j}(t) = \begin{cases}
    0 & g(x) = 0 \quad \ \text{deterministic}\\
    1 & g(x) = x \quad \ \text{demographic}\\
    \mu_{i\setminus j}(t) & g(x) = x^2 \quad \text{environmental}.
  \end{cases}
\end{equation}

An equivalent set of equations can be obtained for the full marginal process:
\begin{align}
  \frac{d}{dt}\mu_{i}(t) & = \langle \dot{x}_{i}(t) \rangle_{i} = \mu_{i}(t)-C_{i}(t,t)-\mu_{i}^2(t)+\sum_{k\in\partial i}J_{ik}\mu_{k\setminus i}(t)\mu_{i}(t) \nonumber \\
  &\quad+ \sum_{k\in \partial i}J_{ik}J_{ki}\int_{0}^{t}ds\, R_{k\setminus i}(t, s)\left(C_{i}(t, s)+\mu_{i}(t)\mu_{i}(s)\right) + \lambda,\\
  \frac{\partial}{\partial t}C_{i}(t,t') & = \langle \dot{x}_{i}(t) x_{i}(t') \rangle_{i} = \left(1 - 2 \mu_{i}(t) \right) C_{i}(t,t') + \sum_{k\in\partial i}J_{ik}\mu_{k\setminus i}(t)C_{i}(t,t') \nonumber\\
  & \quad + \sum_{k\in \partial i}J_{ik}J_{ki}\int_0^t ds\, R_{k\setminus i}(t,s)\left(\mu_{i}(t)C_{i}(s,t') + C_{i}(t,t')\mu_{i}(s)\right) + 2T \mathcal{L}_{i}(t)R_{i}(t',t) \nonumber \\
  & \quad + \sum_{k\in \partial i} J_{ik}^2 \left(\int_0^{t'} ds \, C_{k\setminus i}(t,s) \mu_{i}(t) R_{i}(t',s) + \int_0^{t} ds \, C_{k\setminus i}(t,s) R_{i}(t,s) \mu_{i}(t') \right),\\
  \frac{\partial}{\partial t}R_{i}(t,t') & =  \langle \dot{x}_{i}(t) {\rm i} \hat{x}_{i}(t') \rangle_{i} = \left(1 - 2 \mu_{i}(t) \right) R_{i}(t,t') + \sum_{k\in\partial i}J_{ik}\mu_{k\setminus i}(t)R_{i}(t,t') + \delta(t-t')  \nonumber\\
  & \quad + \sum_{k\in \partial i} J_{ik}J_{ki} \left( \int_{0}^{t} ds \, R_{k\setminus i}(t,s)R_{i}(t,t')\mu_{i}(s)+ \int_{t'}^{t} ds \, \mu_{i}(t)R_{k\setminus i}(t,s) R_{i}(s,t') \right),
\end{align}
where the function $\mathcal{L}_{i}(t)$ is defined as
\begin{equation}
  \mathcal{L}_{i}(t) = \begin{cases}
    0 & g(x) = 0 \quad \ \text{deterministic}\\
    1 & g(x) = x \quad \ \text{demographic}\\
    \mu_{i}(t) & g(x) = x^2 \quad \text{environmental}.
  \end{cases}
\end{equation}

\subsubsection{Stationary state}
Assuming the system reaches a stationary state, we can set $\mu(t)=\mu$, $C(t,t')=C_{i}(t-t')$, and $R(t,t')=R(t-t')$. This corresponds to assuming the system to be in the single equilibrium state. The self-consistent cavity equations then simplify to
\begin{align}
  \frac{d}{dt}\mu_{i\setminus j} = 0 & = \mu_{i\setminus j}\left(1 - \mu_{i\setminus j} + \sum_{k\in\partial i\setminus j}J_{ik}\mu_{k\setminus i} + \sum_{k\in \partial i \setminus j}J_{ik}J_{ki}\chi_{k\setminus i}\mu_{i\setminus j} \right)-C_{i\setminus j}(0)\nonumber \\
  &\quad + \sum_{k\in \partial i\setminus j}J_{ik}J_{ki}\int_0^{\infty}du\, R_{k\setminus i}(u)C_{i\setminus j}(u) + \lambda,\\
  \frac{d}{d\tau}C_{i\setminus j}(\tau) & = C_{i\setminus j}(\tau)\left( 1 - 2\mu_{i\setminus j} + \sum_{k\in\partial i\setminus j}J_{ik}\mu_{k\setminus i} + \sum_{k\in \partial i \setminus j}J_{ik}J_{ki}\chi_{k\setminus i}\mu_{i\setminus j}\right) + 2T\mathcal{L}_{i\setminus j}R_{i\setminus j}(-\tau) \nonumber\\
  & \quad + \mu_{i\setminus j}\sum_{k\in \partial i\setminus j} J_{ik} J_{ki} \int_0^\infty du\, R_{k\setminus i}(u)C_{i\setminus j}(\tau-u) \nonumber \\
  &\quad + \mu_{i\setminus j}\sum_{k\in \partial i\setminus j} J_{ik}^2 \int_0^\infty du \, \left(C_{k\setminus i}(u-\tau) R_{i\setminus j}(u) + C_{k\setminus i}(u) R_{i\setminus j}(u)\right),\\
  \frac{d}{d\tau}R_{i\setminus j}(\tau) & = R_{i\setminus j}(\tau)\left( 1 - 2\mu_{i\setminus j} + \sum_{k\in\partial i\setminus j}J_{ik}\mu_{k\setminus i} + \sum_{k\in \partial i \setminus j}J_{ik}J_{ki}\chi_{k\setminus i}\mu_{i\setminus j}\right) + \delta(\tau) \nonumber\\
  & \quad + \mu_{i\setminus j} \sum_{k\in \partial i\setminus j} J_{ik}J_{ki} \int_0^\tau du \, R_{k\setminus i}(u)R_{i\setminus j}(\tau-u),
\end{align}
where we have introduced the time shift $\tau=t-t'$, and the susceptibility $\chi_{i\setminus j}=\int_0^\infty du\, R_{i\setminus j}(u)$, which is the integrated response of the population $i$ in the cavity graph $\partial i \setminus j$ to a local perturbation. We have defined the function $\mathcal{L}_{i\setminus j}$ as
\begin{equation}
  \mathcal{L}_{i\setminus j} = \begin{cases}
    0 & g(x) = 0 \quad \ \text{deterministic}\\
    1 & g(x) = x \quad \ \text{demographic}\\
    \mu_{i\setminus j} & g(x) = x^2 \quad \text{environmental}.
  \end{cases}
\end{equation}

\subsubsection{Fixed Point + Laplace}

I have the intuition that we have some chances of going ahead exploiting Laplace as in the linear case for $C$ and $R$, but having a fixed point in $\mu$.

Asumming tha $\mu$ reached the fixed point it is easy to show that:

\begin{eqnarray}
  p \hat{R}_{i\setminus j}(p) + \hat{R}_{i\setminus j}(0) &=& \Big[ 1- 2 \mu_{i\setminus j } + \sum_{k\in \partial i \setminus j} J_{ik} \mu_{i\setminus j } \Big] \hat{R}_{i\setminus j}(p) +\nonumber \\
  &+& \mu_{i\setminus j } \hat{R}_{i\setminus j}(p) \sum_{k\in \partial i \setminus j} J_{ik}J_{ki} (\chi_{k\setminus i}  + \hat{R}_{k \setminus i}(p))
\end{eqnarray}

Similarly:

\begin{eqnarray}
  p \hat{C}_{i\setminus j}(p) + \hat{C}_{i\setminus j}(0) &=& \Big[ 1- 2 \mu_{i\setminus j } + \sum_{k\in \partial i \setminus j} J_{ik} \mu_{k\setminus i } \Big] \hat{C}_{i\setminus j}(p) +  \nonumber \\
  &+& \mu_{i\setminus j } \hat{C}_{i\setminus j}(p) \sum_{k\in \partial i \setminus j} J_{ik}J_{ki} (\chi_{k\setminus i}  + \hat{R}_{k \setminus i}(p)) +  \nonumber \\
 & + &\mu_{i\setminus j } \hat{R}_{i\setminus j}(p) \sum_{k\in \partial i \setminus j} J_{ik}^2 \hat{C}_{k\setminus i}(p) + \mu_{i\setminus j } \chi_{i\setminus j} \sum_{k\in \partial i \setminus j} J_{ik}^2 \hat{C}_{k\setminus i}(p)
\end{eqnarray}
where the last term in this expression is {\it dirty}. I assumed that: $\mathcal{L} \Big[ \int_0^t du C_{k\setminus i}(u) R_{k\setminus i}(u) \Big] = \chi_{i\setminus j} \hat{C}_{k\setminus i}(p)$.

Finally, for the fixed point equation:

\begin{equation}
  0 = \mu_{i\setminus j} (1 + \sum_{k\in \partial i \setminus j} J_{ik} \mu_{k\setminus i} ) - \mu_{i\setminus j}^2 (1 - \sum_{k\in \partial i \setminus j} J_{ik}J_{ki} \chi_{k\setminus i} ) - q_{i\setminus j}(1- \sum_{k\in \partial i \setminus j} J_{ik}J_{ki} \chi_{k\setminus i})
\end{equation}
that looks prettier this way:
\begin{equation}
  0 = \mu_{i\setminus j}^2 - \mu_{i\setminus j}\frac{1 + \sum_{k\in \partial i \setminus j} J_{ik} \mu_{k\setminus i}}{1 - \sum_{k\in \partial i \setminus j} J_{ik}J_{ki} \chi_{k\setminus i}} +q_{i\setminus j}
\end{equation}

One can try a bit more and look for fixed point equations also for $C$. I get for $q$:

\begin{equation}
  0 = \Big[ 1 - 2 \mu_{i\setminus j} + \sum_{k\in \partial i \setminus j} J_{ik} \mu_{k\setminus i} + 2\mu_{i\setminus j}  \sum_{k\in \partial i \setminus j} J_{ik}J_{ki} \chi_{k\setminus i} ) \Big] q_{i\setminus j}+ \mu_{i \setminus j} \chi_{i \setminus j}\sum_{k\in \partial i \setminus j} J_{ik}^2 q_{k \setminus i}
\end{equation}

\subsection{Fixed-point equations for the deterministic model}

We focus on the deterministic model without immigration, where $g(x)=0$ and $\lambda=0$. In this case, for certain values of the parameters $m$, $\sigma$ and $\gamma$, we can assume that the system reaches a steady state which does not depend on the initial condition. We further assume that the dynamics reaches a unique fixed-point $\lim_{t\to\infty}x_{i\setminus j}(t)=x_{i\setminus j}^{*}$, and therefore $\lim_{t\to\infty}\zeta_{i\setminus j}(t)=\zeta_{i\setminus j}^{*}$. The fixed point values $x_{i\setminus j}^{*}$, $\zeta_{i\setminus j}^{*}$ are both random variables. It follows that the mean cavity abundance and the cavity correlation become time-independent variables $\lim_{t\to\infty}\mu_{i\setminus j}(t)=\mu_{i\setminus j}^{*}=\langle x_{i\setminus j}^{*}\rangle_{i\setminus j}^{*}$, and $\lim_{t,t'\to\infty}C_{i\setminus j}(t,t')=q_{i\setminus j}^{*}=\langle ( x_{i\setminus j}^{*})^2\rangle_{i\setminus j}^{*} - (\mu_{i\setminus j}^{*})^2$. Within this assumption, the cavity response function is time-translational-invariant, i.e., $R_{i\setminus j}(t,t')$ depends only on the time difference $\tau=t-t'$. Due to causality, $R_{i\setminus j}(\tau)=0$ for $\tau\leq 0$, and we define the susceptibility as $\chi_{i\setminus j}^{*}=\int_{0}^{\infty}d\tau\,R_{i\setminus j}^{*}(\tau)$. 

The fixed point satisfies the self-consistent equation 
\begin{equation}
x_{i\setminus j}^{*}\left[1-x_{i\setminus j}^{*}+\sum_{k\in\partial i \setminus j}J_{ik}\mu_{k \setminus i}^{*}+\sum_{k\in\partial i \setminus j}J_{ik}J_{ki}\chi_{k \setminus i}^{*}x_{i \setminus j}^{*}+\xi_{i \setminus j}^{*}+h_{i \setminus j}\right]=0,\label{eq:FP_eq_cav_edge}
\end{equation}
where $h_{i \setminus j}$ is a constant external field needed to obtain the integrated cavity response, and the noise is Gaussian with zero mean $\langle \xi_{i\setminus j}^*\rangle$ and variance $\langle (\xi_{i\setminus j})^*\rangle=\sum_{k\in\partial i\setminus j}J_{ik}^2q_{k\setminus i}^*=\Sigma_{i\setminus j}^2$. We can therefore write $\xi_{i\setminus j}^*=-\Sigma_{i\setminus j}s$, where $s$ is a standard Gaussian variable with zero mean and unit variance. Eq. \eqref{eq:FP_eq_cav_edge} admits non-vanishing solution
\begin{equation}
x_{i\setminus j}^{*}\left(s\right)=\frac{\Sigma_{i\setminus j}(\Delta_{i\setminus j}-s)+h_{i\setminus j}}{\Gamma_{i\setminus j}}\Theta\left(\frac{\Sigma_{i\setminus j}(\Delta_{i\setminus j}-s)+h_{i\setminus j}}{\Gamma_{i\setminus j}}\right),\label{eq:FP_sol_cavity}
\end{equation}
where $\Theta(x)$ is the Heaviside step function, which ensures positivity of the fixed-point cavity abundance. We have defined $\Delta_{i\setminus j}=(1+\sum_{k\in\partial i \setminus j}J_{ik}\mu_{k \setminus i}^{*})/\Sigma_{i\setminus j}$ and $\Gamma_{i\setminus j}=1-\sum_{k\in\partial i \setminus j}J_{ik}J_{ki}\chi_{k \setminus i}^{*}$.

By averaging over the Gaussian distribution of $s$ we obtain a set of self-consistent equations for the parameters
\begin{subequations}
\begin{align}
\mu_{i\setminus j}^*&=\frac{\Sigma_{i\setminus j}}{\Gamma_{i\setminus j}}\left[\Theta(\Gamma_{i\setminus j})\int_{-\infty}^{\Delta_{i\setminus j}}Ds\,(\Delta_{i\setminus j}-s)+\Theta(-\Gamma_{i\setminus j})\int_{\Delta_{i\setminus j}}^{\infty}Ds\,(\Delta_{i\setminus j}-s)\right]\\
q_{i\setminus j}^*&=\frac{\Sigma^2_{i\setminus j}}{\Gamma^2_{i\setminus j}}\left[\Theta(\Gamma_{i\setminus j})\int_{-\infty}^{\Delta_{i\setminus j}}Ds\,(\Delta_{i\setminus j}-s)^2+\Theta(-\Gamma_{i\setminus j})\int_{\Delta_{i\setminus j}}^{\infty}Ds\,(\Delta_{i\setminus j}-s)^2\right] - \left(\mu_{i\setminus j}^{*}\right)^2\\
\chi_{i\setminus j}^*&=\frac{1}{\Gamma_{i\setminus j}}\left[\Theta(\Gamma_{i\setminus j})\int_{-\infty}^{\Delta_{i\setminus j}}Ds+\Theta(-\Gamma_{i\setminus j})\int_{\Delta_{i\setminus j}}^{\infty}Ds\right].
\end{align}
\end{subequations}
We observe that the integrals can be written in terms of the error function $\erf(x)=\int_0^{x} ds\,e^{-s^2}2/\sqrt{\pi}$. Defining $\varphi(x)=e^{-x^2/2}/\sqrt{2\pi}$ the normal distribution function and $\Phi(x)=[1+\erf(x/\sqrt{2})]/2$ its cumulative distribution function we obtain

\begin{subequations}
\begin{align}
\mu_{i\setminus j}^*&=\frac{\Sigma_{i\setminus j}}{\Gamma_{i\setminus j}}\left[\Theta(\Gamma_{i\setminus j})\left(\Delta_{i\setminus j}\Phi(\Delta_{i\setminus j})+\varphi(\Delta_{i\setminus j}) \right)+\Theta(-\Gamma_{i\setminus j}) \left(\Delta_{i\setminus j}\Phi(-\Delta_{i\setminus j})-\varphi(-\Delta_{i\setminus j}) \right)\right]\\
q_{i\setminus j}^*&=\frac{\Sigma^2_{i\setminus j}}{\Gamma^2_{i\setminus j}}\left[\Theta(\Gamma_{i\setminus j})\left(\left(1+\Delta_{i\setminus j}^2\right)\Phi(\Delta_{i\setminus j})+\Delta_{i\setminus j}\varphi(\Delta_{i\setminus j})\right)\right.\nonumber\\
&\quad \left.+\Theta(\Gamma_{i\setminus j})\left(\left(1+\Delta_{i\setminus j}^2\right)\Phi(-\Delta_{i\setminus j})-\Delta_{i\setminus j}\varphi(-\Delta_{i\setminus j})\right)  \right] - \left(\mu_{i\setminus j}^{*}\right)^2\\
\chi_{i\setminus j}^*&=\frac{1}{\Gamma_{i\setminus j}}\left[\Theta(\Gamma_{i\setminus j})\Phi(\Delta_{i\setminus j})+\Theta(-\Gamma_{i\setminus j})\Phi(-\Delta_{i\setminus j})\right].
\end{align}
\end{subequations}

We can solve these equations by means of a message-passing algorithm, together with the definitions of $\Delta_{i\setminus j}$, $\Sigma_{i\setminus j}$ and $\Gamma_{i\setminus j}$, which are given by
\begin{subequations}
\begin{align}
\Delta_{i\setminus j}&=\frac{1+\sum_{k\in \partial i\setminus j}J_{ik}\mu_{k\setminus i}^*}{\sqrt{\sum_{k\in \partial i\setminus j}J_{ik}^2 q_{k\setminus i}^*}}\\
\Sigma_{i\setminus j}^2&=\sum_{k\in \partial i\setminus j}J_{ik}^2 q_{k\setminus i}^*\\
\Gamma_{i\setminus j}&=1-\sum_{k\in\partial i \setminus j}J_{ik}J_{ki}\chi_{k \setminus i}^{*}.
\end{align}
\end{subequations}


\begin{acknowledgments}
We wish to acknowledge the support of ....
\end{acknowledgments}

\bibliography{GECaM_Ecosystems.bib}

\appendix

\section{\label{sec:effective_cav_proc}Derivation of the effective cavity process}

The calculation is based on the principles of \cite{taraboloGaussianApproximationDynamic2025}.
The dynamical partition function of the Lotka-Volterra system can
be written as

\begin{equation}
  Z=\left\langle \int\mathcal{D}\vec{x}\prod_{i}p_{i}\left(x_{i}\left(0\right)\right)\delta\left(\text{equation of motion}\right)\right\rangle _{\vec{\eta}},
\end{equation}

where $\int\mathcal{D}\vec{x}=\prod_{i}\int\mathcal{D}x$ is a product
of functional integrals, $\langle\dots\rangle_{\vec{\eta}}=\left(\prod_{i}\int\mathcal{D}\eta_{i}\right)\dots$
is the average over the noise realizations, and $\delta(\dots)$ is
a functional Dirac delta function which restrict the integral to the
paths that satisfy the stochastic dynamics. Using the Fourier representation
of the delta function $\delta(x)\propto\int\mathcal{D}\hat{x}\,e^{-{\rm i}\int dt\,\hat{x}(t)x(t)}$
we obtain

\begin{equation}
Z=\left\langle \int\mathcal{D}\vec{x}\mathcal{D}\vec{\hat{x}}\prod_{i}p_{i}\left(x_{i}(0)\right)e^{\int dt\,({-{\rm i}}\hat{x}_{i}(t))\left[\dot{x}_{i}(t)-x_{i}(t)\left(1-x_{i}+\sum_{j\neq i}a_{ij}J_{ij}x_{j}+h_{i}(t)\right)-\eta_i(t)-\lambda\right]}\right\rangle _{\vec{\eta}}.
\end{equation}

The integral over time run over a fixed time interval, but we omit the limits of integration, here and elsewhere, for the sake of simplifying the notation. We introduce an auxiliary variable $f_i=1-x_i-\sum_{j\neq i}a_{ij}J_{ij}x_j(t)+h_i(t)$, which can be inserted in the calculation through a Dirac delta function. Integrating over the Gaussian noise we obtain

\begin{align}
Z =\int\mathcal{D}\vec{x}\mathcal{D}\vec{\hat{x}}\mathcal{D}\vec{f}\mathcal{D}\vec{\hat{f}}\prod_{i\in V}p_{i}\left(x_{i}(0)\right)e^{S_i\left(x_i,\hat{x}_i,f_i,\hat{f}_i\right)}\prod_{(i,j)\in E}e^{S_{ij}\left(x_i,\hat{x}_i,f_i,\hat{f}_i,x_j,\hat{x}_j,f_j,\hat{f}_j\right)},
\end{align}
where we have defined the local and interactive actions as
\begin{align}
  S_i & = \int dt \left\{-{\rm i}\hat{x}_i(t)\left(\dot{x}_i(t)-x_i(t)f_i(t)-\lambda\right)-\frac{1}{2}2Tg(x_i(t))\hat{x}_i^2(t)-{\rm i}\hat{f}_i(t)\left[f_i(t)-\left(1-x_i(t)+h_i(t)\right)\right] \right\}\\
  S_{ij} & = -\int dt\left(J_{ij}{\rm i}\hat{f}_i(t)x_j(t)+J_{ji}{\rm i}\hat{f}_j(t)x_i(t) \right)
\end{align}

We put forward the following \textit{dynamic cavity equations} as an ansatz for describing linearly coupled stochastic dynamics on sparse networks

\begin{equation}
  c_{i\setminus j}[x_{i},\hat{x}_{i}, f_i, \hat{f}_i] = \frac{1}{Z_{i\setminus j}}e^{S_i}\prod_{k\in\partial i \setminus j}\int{\mathcal{D}x_k\mathcal{D}\hat{x}_k\mathcal{D}f_k\mathcal{D}\hat{f}_k} \, c_{k\setminus i}[x_{k},\hat{x}_{k}, f_k, \hat{f}_k] e^{S_{ki}},\label{eq:cavity_marg}
\end{equation}

The full dynamic cavity marginals can be obtained from the cavity
ones as follows 

\begin{equation}
  c_{i}[x_{i},\hat{x}_{i}, f_i, \hat{f}_i]=\frac{1}{Z_{i}}e^{S_i}\prod_{k\in\partial i}\int{\mathcal{D}x_k\mathcal{D}\hat{x}_k\mathcal{D}f_k\mathcal{D}\hat{f}_k} \, c_{k\setminus i}[x_{k},\hat{x}_{k}, f_k, \hat{f}_k] e^{S_{ki}}.\label{eq:marg_RRG}
\end{equation}

Expanding the cavity marginals Eq. \eqref{eq:cavity_marg} at
second order in the coupling parameters $J_{ij}$ we obtain an effective single node action for the cavity marginal

\begin{equation}
  c_{i\setminus j}[x_{i},\hat{x}_{i}, f_i, \hat{f}_i] \approx \frac{p_i(x_i(0))}{Z_{i\setminus j}}e^{S^{\rm eff}_{i\setminus j}},
\end{equation}

where the effective cavity action is defined as

\begin{align}
  S^{\rm eff}_{i\setminus j} & = \int dt \Bigg\{-{\rm i}\hat{x}_i(t)\left(\dot{x}_i(t)-x_i(t)f_i(t)-\lambda\right)-\frac{1}{2}2Tg(x_i(t))\hat{x}_i^2(t)-{\rm i}\hat{f}_i(t)\Big[f_i(t)-\left(1-x_i(t)+h_i(t)\right) \nonumber \\
  & \quad  + \sum_{k\in\partial i \setminus j}J_{ik}\mu_{k\setminus i}(t)- \sum_{k\in\partial i \setminus j}J_{ik}J_{ki}\int dt' \, R_{k\setminus i}(t,t')x_i(t')\Big]\Bigg\}+\int dtdt' \, {\rm i}\hat{f}(t)\left(  \sum_{k\in\partial i \setminus j}J_{ik}^2 C_{k\setminus i}(t,t')\right){\rm i}\hat{f}(t')
\end{align}

where $\mu_{k\setminus i}(t)$, $R_{k\setminus i}(t,t')$, and $C_{k\setminus i}(t,t')$
are respectively the cavity mean abundance, the cavity response function, and the cavity (connected) two-point correlation function. They are obtained self-consistently from the effective dynamics of the neighbors

\begin{equation}
\mu_{k\setminus i}(t)=\left\langle x_{k}(t)\right\rangle _{k\setminus i},\quad R_{k\setminus i}(t,t')=\left\langle x_{k}(t){\rm i}\hat{f}_{k}(t')\right\rangle _{k\setminus i},\quad C(t,t')=\left\langle x_{k}(t)x_{k}(t')\right\rangle _{k\setminus i} - \mu_{k\setminus i}(t)\mu_{k\setminus i}(t'),
\end{equation}

where $\left\langle \dots\right\rangle _{k\setminus i}=\int{\mathcal{D}x_k\mathcal{D}\hat{x}_k\mathcal{D}f_k\mathcal{D}\hat{f}_k} \, c_{k\setminus i}[x_{k},\hat{x}_{k}, f_k, \hat{f}_k] \dots$
describes the local average over the cavity marginal.

The effective cavity action can be identified with a sort of local effective single-edge action for the variables on site $i$ when the interaction term with site $j$ is removed from the underlying graph.

Performing backward the steps of the path-integral representation, we obtain an effective stochastic dynamics in terms of the stochastic differential equation on the cavity graph

\begin{align}
  \frac{d}{dt}x_{i\setminus j}(t) & =x_{i\setminus j}(t)\Biggl[1-x_{i\setminus j}(t)+\sum_{k\in\partial i\setminus j}J_{ik}\mu_{k\setminus i}(t) + \sum_{k\in\partial i\setminus j}\int_{0}^{t}dt'J_{ik}J_{ki}R_{k\setminus i}(t,t')x_{i\setminus j}(t') \nonumber \\
  & \qquad\quad \qquad + h_{i\setminus j}(t)+\zeta_{i\setminus j}(t)\Biggr]+\sqrt{2Tg(x_{i\setminus j})}\xi_{i\setminus j}(t)+\lambda,
\end{align}

where $\zeta_{i\setminus j}$ and $\xi_{i\setminus j}$ are Gaussian noises with vanishing mean and correlations

\begin{equation}
  \left\langle \zeta_{i\setminus j}(t)\zeta_{i\setminus j}(t')\right\rangle _{i\setminus j}=\sum_{k\in\partial i\setminus j}J_{ik}^{2}C_{k\setminus i}(t,t'),\qquad\left\langle \xi_{i\setminus j}(t)\xi_{i\setminus j}(t')\right\rangle _{i\setminus j}=\delta(t-t').
\end{equation}

The terms $\mu_{k\setminus i}(t)$, $R_{k\setminus i}(t,t')$, and $C_{k\setminus i}(t,t')$ have to be considered as external parameters, obtained self-consistently as averages over the effective dynamics of the neighbors. In a similar way, an effective process for the original graph can be obtained as

\begin{align}
  \frac{d}{dt}x_{i}(t) & =x_{i}(t)\Biggl[1-x_{i}(t)+\sum_{k\in\partial i}J_{ik}\mu_{k\setminus i}(t) + \sum_{k\in\partial i}\int_{0}^{t}dt'J_{ik}J_{ki}R_{k\setminus i}(t,t')x_{i}(t') \nonumber \\
   & \qquad\quad \qquad + h_{i}(t)+\zeta_{i}(t)\Biggr]+\sqrt{2Tg(x_{i}(t))}\xi_{i}(t)+\lambda,
\end{align}

with Gaussian noises with vanishing mean and correlations

\begin{equation}
  \left\langle \zeta_{i}(t)\zeta_{i}(t')\right\rangle _{i}=\sum_{k\in\partial i}J_{ik}^{2}C_{k\setminus i}(t,t'),\qquad\left\langle \xi_{i}(t)\xi_{i}(t')\right\rangle _{i}=\delta(t-t').
\end{equation}


\section{Single-path response function}\label{app:response}
The single-path response function $R^{\rm sp}_{i\setminus j}(t,t')=\delta/\delta h_{i\setminus j}(t') x_{i\setminus j}(t) $ can be obtained by directly applying the functional derivative $\delta/\delta h_{i\setminus j}(t')$ to Eq.~\eqref{eq:eff_proc_cav}:
\begin{align}
  \frac{\partial}{\partial t}R^{\rm sp}_{i\setminus j}(t,t')&=R^{\rm sp}_{i\setminus j}(t,t')\left(1-2x_{i\setminus j}(t)\right)+R^{\rm sp}_{i\setminus j}(t,t')\sqrt{\frac{2T}{g(x_{i\setminus j}(t))}}g'(x_{i\setminus j}(t))\xi_{i\setminus j}(t)\nonumber\\
  &\quad + R^{\rm sp}_{i\setminus j}(t,t')\left(\sum_{k\in\partial i \setminus j}J_{ik}\mu_{k\setminus i}(t)+\zeta_{i\setminus j}(t) \right)+x_{i\setminus j}(t)\delta(t-t')\nonumber\\
  &\quad + R^{\rm sp}_{i\setminus j}(t,t')\sum_{k \in \partial i \setminus j}J_{ik}J_{ki}\int_0^t dsR_{k\setminus i}(t,s)x_{i\setminus j}(s)+x_{i\setminus j}(t)\sum_{k\in \partial i \setminus j}J_{ik}J_{ki}\int_0^t ds R_{k\setminus i}(t,s)R^{\rm sp}_{i\setminus j}(s,t'). \label{eq:single_path_response}
\end{align}
The single-path response function can therefore be integrated numerically using the initial condition $R^{\rm sp}_{i\setminus j}(t,t')=0$ for $t<t'$ from causality.

\section{Numerical integration of GECaM dynamics}\label{app:numerical_integration}

The numerical integration of the GECaM dynamics is performed using an Euler-Maruyama scheme with a time step $\Delta$. Time is discretized in intervals of size $\Delta$, i.e. $t=n\Delta$ with $n=0,1,\ldots,T$, where $T$ is the number of time steps. The trajectories $x(t)$ become $T$-dimensional time-vectors with elements $x^n = x(t=n\Delta)$, while the self-consistent fields $\mu(t)$, $R(t,t')$ and $C(t,t')$ become, respectively, $T$-dimensional vectors and $T\times T$ matrices, with elements $\mu^n = \mu(t=n\Delta)$, $R^{n,l}=R(t=n\Delta,t'=l\Delta)$ and $C^{n,l}=C(t=n\Delta,t'=l\Delta)$.

Both the effective cavity SDEs Eq.~\eqref{eq:eff_proc_cav} and the effective marginal SDEs Eq.~\eqref{eq:eff_proc_marg} are better integrated in $\log$ space. The discretized effective cavity SDE in $\log$ space reads
\begin{equation}\label{eq:log_eff_cav_discrete}
  \log x_{i\setminus j}^{n+1} = \log x_{i\setminus j}^n + \Delta \mathcal{F}^n_{i\setminus j}\left(\bm{x}_{i\setminus j}\Big|\left\{\mu^n_{k\setminus i}\right\}_{k\in \partial i\setminus j},\left\{\bm{R}_{k\setminus i}\right\}_{k\in \partial i\setminus j},\zeta_{i\setminus j}^n\right) + \Delta\sqrt{2T}\mathcal{G}\left(x_{i\setminus j}^n\big|\xi_{i\setminus j}^n\right) + \Delta \mathcal{H}\left(x_{i\setminus j}^n\big|\lambda\right),
\end{equation}
where
\begin{align}
  \mathcal{F}^n_{i\setminus j}\left(\bm{x}_{i\setminus j}\Big|\left\{\mu^n_{k\setminus i}\right\}_{k\in \partial i\setminus j},\left\{\bm{R}_{k\setminus i}\right\}_{k\in \partial i\setminus j},\zeta_{i\setminus j}^n\right) & = 1-x_{i\setminus j}^n+\sum_{k\in \partial i \setminus j}J_{ik}\mu^n_{k\setminus i} + \Delta \sum_{m=0}^{n-1} x_{i\setminus j}^{m} \sum_{k \in \partial i \setminus j} J_{ik}J_{ki}R_{k\setminus i}^{n,m} + \zeta_{i\setminus j}^n,\\
  \mathcal{G}\left(x\mid\xi\right) & = \begin{cases}
    0 & g(x) = 0 \quad \ \text{deterministic}\\
    \exp\left[\log \xi - \frac{1}{2}\log x \right] & g(x) = x \quad \ \text{demographic}\\
    \xi & g(x) = x^2 \quad \text{environmental},
  \end{cases}\\
  \mathcal{H}\left(x\mid \lambda\right) & = \begin{cases}
    0 & \lambda = 0 \\
    \exp\left[\log \lambda - \log x\right] & \lambda > 0.
  \end{cases}
\end{align}

The discretized equation for the single path cavity response function reads
\begin{align}
  R^{({\rm sp})n+1,l}_{i\setminus j}&=R^{({\rm sp})n,l}_{i\setminus j}+\Delta R^{({\rm sp})n,l}_{i\setminus j}\left(1-2x^n_{i\setminus j} + \sum_{k\in\partial i \setminus j}J_{ik}\mu_{k\setminus i}^n+\zeta^n_{i\setminus j} + \sqrt{2T}\xi_{i\setminus j}^n\mathcal{I}\left(x_{i\setminus j}^n\right) \right)+x^n_{i\setminus j}\delta_{n,l}\nonumber\\
  &\qquad + \Delta^2 R^{(\rm sp)n,l}_{i\setminus j}\sum_{m=0}^{n-1}x_{i\setminus j}^m\sum_{k \in \partial i \setminus j}J_{ik}J_{ki} R_{k\setminus i}^{n,m} + \Delta^2 x_{i\setminus j}^n\sum_{m=l+1}^{n-1}R^{(\rm sp)m,l}_{i\setminus j}\sum_{k\in \partial i \setminus j}J_{ik}J_{ki} R_{k\setminus i}^{n,m}, \label{eq:single_path_response_discrete}
\end{align}
where
\begin{align}
  \mathcal{I}\left(x\right) & = \begin{cases}
    0 & g(x) = 0 \quad \ \text{deterministic}\\
    1/\sqrt{x} & g(x) = x \quad \ \text{demographic}\\
    1 & g(x) = x^2 \quad \text{environmental}.
  \end{cases}
\end{align}
The discretized equation for the single path marginal response function reads
\begin{align}
  R^{({\rm sp})n+1,l}_{i}&=R^{({\rm sp})n,l}_{i}+\Delta R^{({\rm sp})n,l}_{i}\left(1-2x^n_{i} + \sum_{k\in\partial i}J_{ik}\mu_{k\setminus i}^n+\zeta^n_{i} + \sqrt{2T}\xi_{i}^n\mathcal{I}\left(x_{i}^n\right) \right)+x^n_{i}\delta_{n,l}\nonumber\\
  &\qquad + \Delta^2 R^{(\rm sp)n,l}_{i}\sum_{m=0}^{n-1}x_{i}^m\sum_{k \in \partial i}J_{ik}J_{ki} R_{k\setminus i}^{n,m} + \Delta^2 x_{i}^n\sum_{m=l+1}^{n-1}R^{(\rm sp)m,l}_{i}\sum_{k\in \partial i}J_{ik}J_{ki} R_{k\setminus i}^{n,m}. \label{eq:single_path_response_marg_discrete}
\end{align}

The discretized effective marginal SDE in $\log$ space reads
\begin{equation}\label{eq:log_eff_marg_discrete}
  \log x_{i}^{n+1} = \log x_{i}^n + \Delta \mathcal{F}^n_{i}\left(\bm{x}_{i}\Big|\left\{\mu^n_{k\setminus i}\right\}_{k\in \partial i},\left\{\bm{R}_{k\setminus i}\right\}_{k\in \partial i},\zeta_{i}^n\right) + \Delta\sqrt{2T}\mathcal{G}\left(x_{i}^n\big|\xi_{i}^n\right) + \Delta \mathcal{H}\left(x_{i}^n\big|\lambda\right),
\end{equation}
where
\begin{align}
  \mathcal{F}^n_{i}\left(\bm{x}_{i}\Big|\left\{\mu^n_{k\setminus i}\right\}_{k\in \partial i},\left\{\bm{R}_{k\setminus i}\right\}_{k\in \partial i},\zeta_{i}^n\right) & = 1-x_{i}^n+\sum_{k\in \partial i}J_{ik}\mu^n_{k\setminus i} + \Delta \sum_{m=0}^{n-1} x_{i}^{m} \sum_{k \in \partial i} J_{ik}J_{ki}R_{k\setminus i}^{n,m} + \zeta_{i}^n.
\end{align}

In order to avoid unnecessary computations, we can define the following node quantities:
\begin{align}
  \Sigma_{i,\mu}^{n} & = \sum_{k\in \partial i}J_{ik}\mu_{k\setminus i}^n,\label{eq:sum_mu}\\
  \Sigma_{i,R}^{n,l} & = \sum_{k\in \partial i}J_{ik}J_{ki} R_{k\setminus i}^{n,l},\label{eq:sum_R}\\
  \Sigma_{i,C}^{n,l} & = \sum_{k\in \partial i}J_{ik}^2 C_{k\setminus i}^{n,l},\label{eq:sum_C}\\
\end{align}
which will be computed once at the beginning of the integration. The cavity contributions can therefore be computed by subtracting the contribution of the node $j$ from the sum over the neighbors of node $i$, thus avoiding the need to compute the sum over all neighbors at each time step.

The resulting algorithm proceeds as follows. 

\begin{enumerate}
  \item Propose initial guesses for the cavity mean abundances $\{\mu_{i\setminus j}\}_{(i,j)\in E}$, the cavity correlations $\{C_{i\setminus j}\}_{(i,j)\in E}$, and the cavity responses $\{R_{i\setminus j}\}_{(i,j)\in E}$ for each edge of the graph.
  \item For each node $i$ of the graph:
  \begin{enumerate} 
    \item Compute the node quantities $\Sigma_{i,\mu}^n$, $\Sigma_{i,R}^{n,l}$ and $\Sigma_{i,C}^{n,l}$ using Eqs. \eqref{eq:sum_mu}, \eqref{eq:sum_R} and \eqref{eq:sum_C}.
    \item For each edge $(i,j)$ of the graph:
    \begin{enumerate}
      \item Sample $N_{\rm traj}$ Gaussian paths $\{(\bm{\zeta}^{(s)}_{i\setminus j},\bm{\xi}^{(s)}_{i\setminus j})\}_{s=1}^{N_{\rm traj}}$, where $\bm{\zeta}_{i\setminus j}$ and $\bm{\xi}_{i\setminus j}$ have, respectively, covariance matrices $C^{(\zeta)n,l}_{i\setminus j}=\Sigma_{i,C}^{n,l} - J_{ik}^2C^{n,l}_{k\setminus i}$ and $C^{(\xi)n,l}_{i\setminus j}=\delta_{n,l}$. This can be done by diagonalizing the covariance matrix $C^{(\zeta)}_{i\setminus j}$ and using the proper basis to sample independent gaussian components.
      \item For each path $(\bm{\zeta}^{(s)}_{i\setminus j},\bm{\xi}^{(s)}_{i\setminus j})$, numerically integrate the effective cavity equation with an initial condition sampled according to $p_0(x)$ using the Euler-Maruyama scheme in $\log$ space, i.e. Eq. \eqref{eq:log_eff_cav_discrete}. This will yield $N_{\rm traj}$ trajectories $\{\bm{x}_{i\setminus j}^{(s)}\}_{s=1}^{N_{\rm traj}}$.
      \item For each path $(\bm{\zeta}^{(s)}_{i\setminus j},\bm{\xi}^{(s)}_{i\setminus j})$, compute the single-path response function using Eq. \eqref{eq:single_path_response_discrete}. This will yield $N_{\rm traj}$ single path response matrices $\{\bm{R}^{(s)}_{i\setminus j}\}_{s=1}^{N_{\rm traj}}$.
      \item Compute the new values of the cavity self-consistent fields as averages over the paths:
      \begin{align*}
        \mu_{i\setminus j}^{({\rm new})n} & = \frac{1}{N_{\rm traj}}\sum_{s=1}^{N_{\rm traj}}x_{i\setminus j}^{(s)n},\\
        C_{i\setminus j}^{({\rm new})n,l} & = \frac{1}{N_{\rm traj}}\sum_{s=1}^{N_{\rm traj}}x_{i\setminus j}^{(s)n} x_{i\setminus j}^{(s)l} - \mu_{i\setminus j}^n \mu_{i\setminus j}^l,\\
        R_{i\setminus j}^{({\rm new})n,l} & = \frac{1}{N_{\rm traj}}\sum_{s=1}^{N_{\rm traj}}R^{(s)n,l}_{i\setminus j,k}.
      \end{align*}
    \end{enumerate}
  \end{enumerate}
  \item Check for convergence of the cavity self-consistent fields. 
  \begin{equation}
    \max_{(i,j)}\left\lVert C^{\rm updated}_{i\setminus j} - C_{i\setminus j}^{\rm new} \right\rVert < \epsilon,
  \end{equation}
  If not converged, go to step 2.
  \item For each node $i$ of the graph:
  \begin{enumerate}
    \item Compute the node quantities $\Sigma_{i,\mu}^n$, $\Sigma_{i,R}^{n,l}$ and $\Sigma_{i,C}^{n,l}$ using Eqs. \eqref{eq:sum_mu}, \eqref{eq:sum_R} and \eqref{eq:sum_C}.
    \item Sample $N_{\rm traj}$ Gaussian paths $\{(\bm{\zeta}^{(s)}_{i},\bm{\xi}^{(s)}_{i})\}_{s=1}^{N_{\rm traj}}$, where $\bm{\zeta}_{i}$ and $\bm{\xi}_{i}$ have, respectively, covariance matrices $C^{(\zeta)n,l}_{i}=\Sigma_{i,C}^{n,l}$ and $C^{(\xi)n,l}_{i}=\delta_{n,l}$. This can be done by diagonalizing the covariance matrix $C^{(\zeta)}_{i}$ and using the proper basis to sample independent gaussian components.
    \item For each path $(\bm{\zeta}^{(s)}_{i},\bm{\xi}^{(s)}_{i})$, numerically integrate the effective marginal equation with an initial condition sampled according to $p_0(x)$ using the Euler-Maruyama scheme in $\log$ space, i.e. Eq. \eqref{eq:log_eff_marg_discrete}. This will yield $N_{\rm traj}$ trajectories $\{\bm{x}_{i}^{(s)}\}_{s=1}^{N_{\rm traj}}$.
    \item For each path $(\bm{\zeta}^{(s)}_{i},\bm{\xi}^{(s)}_{i})$, compute the single-path response function using Eq. \eqref{eq:single_path_response_discrete}. This will yield $N_{\rm traj}$ single path response matrices $\{\bm{R}^{(s)}_{i}\}_{s=1}^{N_{\rm traj}}$.
    \item Compute the marginal self-consistent fields as averages over the paths:
    \begin{align*}
      \mu_{i}^{n} & = \frac{1}{N_{\rm traj}}\sum_{s=1}^{N_{\rm traj}}x_{i}^{(s)n},\\
      C_{i}^{n,l} & = \frac{1}{N_{\rm traj}}\sum_{s=1}^{N_{\rm traj}}x_{i}^{(s)n} x_{i}^{(s)l} - \mu_{i}^n \mu_{i}^l,\\
      R_{i}^{n,l} & = \frac{1}{N_{\rm traj}}\sum_{s=1}^{N_{\rm traj}}R^{(s)n,l}_{i,k}.
    \end{align*}
  \end{enumerate}
\end{enumerate}

\section{Gaussian closure}
The effective process Eq. \eqref{eq:eff_proc_marg} is a stochastic differential equation that can be solved numerically. However, it is often useful to obtain analytical closed form equations for the fields $\mu$, $C$ and $R$. 
For simplicity, we rewrite the effective cavity process in the form
\begin{equation}
  \frac{d}{dt}x_{i\setminus j}(t) = x_{i\setminus j}(t)\Biggl[1-x_{i\setminus j}(t)+\sum_{k\in\partial i\setminus j}J_{ik}\mu_{k\setminus i}(t) + \sum_{k\in\partial i\setminus j}\int_{0}^{t}dt'J_{ik}J_{ki}R_{k\setminus i}(t,t')x_{i\setminus j}(t')\Biggr] + \eta_{i\setminus j}(t)+\lambda,
\end{equation}
where we have defined the noise term $\eta_{i\setminus j}(t)$ as a Gaussian process with vanishing mean and covariance
\begin{equation}
  \left\langle \eta_{i\setminus j}(t)\eta_{i\setminus j}(t')\right\rangle = x_{i\setminus j}(t)\sum_{k\in\partial i\setminus j}J_{ik}^{2}C_{k\setminus i}(t,t')x_{i\setminus j}(t') + 2Tg(x_{i\setminus j}(t))\delta(t-t').
\end{equation}
The approximated cavity message can be written as
\begin{subequations}
  \begin{align}
    c_{i\setminus j}[x_{i\setminus j},\hat{x}_{i\setminus j}] &= \Bigg\langle\exp \int dt \Bigg\{-{\rm i}\hat{x}_{i\setminus j}(t)\Bigg[\frac{\dot{x}_{i\setminus j}(t)}{x_{i\setminus j}(t)}-\Big(1-x_{i\setminus j}(t)+\sum_{k\in\partial i\setminus j}J_{ik}\mu_{k\setminus i}(t)\nonumber \\
    &\quad + \sum_{k\in\partial i\setminus j}\int_{0}^{t}dsJ_{ik}J_{ki}R_{k\setminus i}(t,s)x_{i\setminus j}(s) + \frac{\eta_{i\setminus j}(t)}{x_{i\setminus j}(t)} + \frac{\lambda}{x_{i\setminus j}(t)} \Big)\Bigg]\Bigg\}\Bigg\rangle\\
    &= \exp \Bigg\{-\int dt\, \frac{{\rm i}\hat{x}_{i\setminus j}(t)}{x_{i\setminus j}(t)}\Bigg[\dot{x}_{i\setminus j}(t)-x_{i\setminus j}(t)\Big(1-x_{i\setminus j}(t)+\sum_{k\in\partial i\setminus j}J_{ik}\mu_{k\setminus i}(t)\nonumber \\
    &\quad + \sum_{k\in\partial i\setminus j}\int_{0}^{t}dsJ_{ik}J_{ki}R_{k\setminus i}(t,s)x_{i\setminus j}(s)\Big) - \lambda\Bigg] \Bigg\} \left \langle \exp \int dt \frac{{\rm i}\hat{x}_{i\setminus j}(t)}{x_{i\setminus j}(t)} \eta_{i\setminus j}(t)\right \rangle\\
    &= \exp \Bigg\{-\int dt\, \frac{{\rm i}\hat{x}_{i\setminus j}(t)}{x_{i\setminus j}(t)}\Bigg[\dot{x}_{i\setminus j}(t)-x_{i\setminus j}(t)\Big(1-x_{i\setminus j}(t)+\sum_{k\in\partial i\setminus j}J_{ik}\mu_{k\setminus i}(t)\nonumber \\
    &\quad + \sum_{k\in\partial i\setminus j}\int_{0}^{t}dsJ_{ik}J_{ki}R_{k\setminus i}(t,s)x_{i\setminus j}(s)\Big) - \lambda\Bigg] \nonumber\\
    &\quad + \frac{1}{2} \int dt dt' \, \frac{{\rm i}\hat{x}_{i\setminus j}(t)}{x_{i\setminus j}(t)}\left(2T g(x_{i\setminus j}(t))\delta(t-t')+x_{i\setminus j}(t)\sum_{k\in \partial i \setminus j}J_{ik}^2 C_{k\setminus i}(t,t') x_{i\setminus j}(t')\right)\frac{{\rm i}\hat{x}_{i\setminus j}(t')}{x_{i\setminus j}(t')}\Bigg\}.
  \end{align}
\end{subequations}

By averaging Eq. \eqref{eq:eff_proc_cav} over the noise realizations, we obtain the following self-consistent equations for the cavity mean abundances
\begin{subequations}
  \begin{align}
    \frac{d}{dt}\mu_{i\setminus j}(t) & =  \langle \dot{x}_{i\setminus j}(t) \rangle = \mu_{i\setminus j}(t)-C_{i\setminus j}(t,t)-\mu_{i\setminus j}^2(t)+\sum_{k\in\partial i \setminus j}J_{ik}\mu_{k\setminus i}(t)\mu_{i\setminus j}(t)  \nonumber \\
    &\quad +\sum_{k\in \partial i \setminus j}J_{ik}J_{ki}\int_{0}^{t}ds R_{k\setminus i}(t, s)\left(C_{i\setminus j}(t, s)+\mu_{i\setminus j}(t)\mu_{i\setminus j}(s)\right)+ \left\langle \eta_{i\setminus j}(t)\right\rangle + \lambda\\
    & = \mu_{i\setminus j}(t)-C_{i\setminus j}(t,t)-\mu_{i\setminus j}^2(t)+\sum_{k\in\partial i \setminus j}J_{ik}\mu_{k\setminus i}(t)\mu_{i\setminus j}(t)  \nonumber \\
    &\quad +\sum_{k\in \partial i \setminus j}J_{ik}J_{ki}\int_{0}^{t}ds R_{k\setminus i}(t, s)\left(C_{i\setminus j}(t, s)+\mu_{i\setminus j}(t)\mu_{i\setminus j}(s)\right) + \lambda.
  \end{align}
\end{subequations}
Similarly, we can obtain the self-consistent equations for the disconnected cavity correlations
\begin{subequations}
  \begin{align}
    \frac{\partial}{\partial t}C^{\rm dc}_{i\setminus j}(t,t') & = \langle \dot{x}_{i\setminus j}(t) x_{i\setminus j}(t') \rangle = C^{\rm dc}_{i\setminus j}(t,t') - \langle x_{i\setminus j}^2(t)x_{i\setminus j}(t') \rangle + \sum_{k\in\partial i \setminus j}J_{ik}\mu_{k\setminus i}(t)C^{\rm dc}_{i\setminus j}(t,t') \nonumber\\
    &\quad + \sum_{k\in \partial i \setminus j}J_{ik}J_{ki}\int_0^t ds R_{k\setminus i}(t,s)\langle x_{i\setminus j}(t)x_{i\setminus j}(s)x_{i\setminus j}(t')\rangle + \left \langle \eta_{i\setminus j}(t)x_{i\setminus j}(t')\right \rangle + \lambda \mu_{i\setminus j}( t')\\
    & = C^{\rm dc}_{i\setminus j}(t,t') - \langle x_{i\setminus j}^2(t)x_{i\setminus j}(t') \rangle + \sum_{k\in\partial i \setminus j}J_{ik}\mu_{k\setminus i}(t)C^{\rm dc}_{i\setminus j}(t,t') \nonumber\\
    &\quad + \sum_{k\in \partial i \setminus j}J_{ik}J_{ki}\int_0^t ds R_{k\setminus i}(t,s)\langle x_{i\setminus j}(t)x_{i\setminus j}(s)x_{i\setminus j}(t')\rangle + 2T \langle \frac{g(x_{i\setminus j}(t))}{x_{i\setminus j}(t)}{\rm i}\hat{x}_{i\setminus j}(t)x_{i\setminus j}(t') \rangle \nonumber \\
    &\quad + \sum_{k\in \partial i \setminus j} J_{ik}^2 \int ds C_{k\setminus i}(t,s) \langle x_{i \setminus j}(t) {\rm i} \hat{x}_{i \setminus j}(s) x_{i \setminus j}(t') \rangle + \lambda \mu_{i\setminus j}( t').
  \end{align}
\end{subequations}
The self-consistent equations for the cavity response functions can be obtained as
\begin{subequations}
  \begin{align}
    \frac{\partial}{\partial t}R_{i\setminus j}(t,t') & = \langle \dot{x}_{i\setminus j}(t){\rm i}\hat{x}_{i\setminus j}(t')\rangle = R_{i\setminus j}(t,t') - \langle x_{i\setminus j}^2(t){\rm i}\hat{x}_{i\setminus j}(t') \rangle + \sum_{k\in\partial i \setminus j}J_{ik}\mu_{k\setminus i}(t)R_{i\setminus j}(t,t') \nonumber\\
    &\quad + \sum_{k\in \partial i \setminus j} J_{ik}J_{ki} \int ds R_{k\setminus i}(t,s) \langle x_{i\setminus j}(t) x_{i\setminus j}(s) {\rm i}\hat{x}_{i\setminus j}(t') \rangle + \langle\eta_{i\setminus j}(t){\rm i}\hat{x}_{i\setminus j}(t') \rangle.\\
    & = R_{i\setminus j}(t,t') - \langle x_{i\setminus j}^2(t){\rm i}\hat{x}_{i\setminus j}(t') \rangle + \sum_{k\in\partial i \setminus j}J_{ik}\mu_{k\setminus i}(t)R_{i\setminus j}(t,t') \nonumber\\
    &\quad + \sum_{k\in \partial i \setminus j} J_{ik}J_{ki} \int ds R_{k\setminus i}(t,s) \langle x_{i\setminus j}(t) x_{i\setminus j}(s) {\rm i}\hat{x}_{i\setminus j}(t') \rangle + 2T \langle \frac{g(x_{i\setminus j}(t))}{x_{i\setminus j}(t)}{\rm i}\hat{x}_{i\setminus j}(t){\rm i} \hat{x}_{i\setminus j}(t') \rangle \nonumber\\
    &\quad + \sum_{k\in \partial i \setminus j} J_{ik}^2 \int ds C_{k\setminus i}(t,s) \langle x_{i \setminus j}(t) {\rm i} \hat{x}_{i \setminus j}(s) {\rm i} \hat{x}_{i \setminus j}(t') \rangle.
  \end{align}
\end{subequations}
The equations for the correlations and responses are not closed, as they depend on three-, four- and five-times averages. We can close the equations by approximating those averages with sums of products of two-times averages, i.e. we can use the Gaussian closure approximation.
Specifically, introducing the zero-mean degrees of freedom $\delta x(t) = x(t) - \mu(t)$, we can approximate as:
\begin{align*}
  \langle x^2(t)x(t') \rangle & = \langle (\delta x(t)+\mu(t))^2 (\delta x(t')+\mu(t')) \rangle \approx C(t,t)\mu(t') + \mu^2(t)\mu(t') + 2\mu(t)C(t,t'),\\
  \langle x(t)x(s)x(t') \rangle & = \langle (\delta x(t)+\mu(t)) (\delta x(s)+\mu(s)) (\delta x(t')+\mu(t')) \rangle \\
  & \approx C(t,s)\mu(t')+C(t,t')\mu(s)+C(s,t')\mu(t)+\mu(t)\mu(s)\mu(t'),\\
  \langle x(t){\rm i}\hat{x}(s)x(t') \rangle & = \langle (\delta x(t)+\mu(t)) {\rm i}\hat{x}(s) (\delta x(t')+\mu(t')) \rangle \approx \mu(t)R(t',s) + R(t,s)\mu(t'),\\
  \langle x^2(t){\rm i}\hat{x}(t') \rangle & = \langle (\delta x(t)+\mu(t))^2 {\rm i}\hat{x}(t') \rangle \approx 2\mu(t)R(t,t'),\\
  \langle x(t)x(s){\rm i}\hat{x}(t') \rangle & = \langle (\delta x(t)+\mu(t)) (\delta x(s)+\mu(s)) {\rm i}\hat{x}(t') \rangle \approx R(t,t')\mu(s)+\mu(t) R(s,t'),\\
  \langle x(t){\rm i}\hat{x}(s){\rm i}\hat{x}(t') \rangle & = \langle (\delta x(t)+\mu(t)) {\rm i}\hat{x}(s) {\rm i}\hat{x}(t') \rangle \approx 0.
\end{align*}
The terms involving the function $g(x)$ depend on the specific form of the function.
\begin{itemize}
\item \textbf{Demographic:} $g(x) = x$
  \begin{align*}
    \langle \frac{g(x(t))}{x(t)}{\rm i}\hat{x}(t)x(t') \rangle & = \langle {\rm i}\hat{x}(t) x(t')\rangle = R(t',t),\\
    \langle \frac{g(x(t))}{x(t)}{\rm i}\hat{x}(t){\rm i}\hat{x}(t') \rangle & = \langle {\rm i}\hat{x}(t){\rm i}\hat{x}(t') \rangle = 0.
  \end{align*}
\item \textbf{Environmental:} $g(x) = x^2$
  \begin{align*}
    \langle \frac{g(x(t))}{x(t)}{\rm i}\hat{x}(t)x(t') \rangle & = \langle (\delta x(t)+\mu(t)){\rm i}\hat{x}(t)(\delta x(t')+\mu(t')) \rangle \approx \mu(t)R(t',t),\\
    \langle \frac{g(x(t))}{x(t)}{\rm i}\hat{x}(t){\rm i}\hat{x}(t') \rangle & = \langle (\delta x(t)+\mu(t)){\rm i}\hat{x}(t){\rm i}\hat{x}(t') \rangle \approx 0.
  \end{align*}
\end{itemize}

The resulting closed equations for the cavity mean abundances, correlations and responses read
\begin{align}
  \frac{d}{dt}\mu_{i\setminus j}(t) & = \mu_{i\setminus j}(t)-C_{i\setminus j}(t,t)-\mu_{i\setminus j}^2(t)+\sum_{k\in\partial i \setminus j}J_{ik}\mu_{k\setminus i}(t)\mu_{i\setminus j}(t)  \nonumber \\
  & \quad +\sum_{k\in \partial i \setminus j}J_{ik}J_{ki}\int_{0}^{t}ds R_{k\setminus i}(t, s)\left(C_{i\setminus j}(t, s)+\mu_{i\setminus j}(t)\mu_{i\setminus j}(s)\right) + \lambda,\\
  \frac{\partial}{\partial t}C^{\rm dc}_{i\setminus j}(t,t') & = C^{\rm dc}_{i\setminus j}(t,t')\left(1-2\mu_{i\setminus j}(t)\right) - C^{\rm dc}_{i\setminus j}(t,t)\mu_{i\setminus j}(t') + 2 \mu_{i\setminus j}^2(t)\mu_{i\setminus j}(t') + \sum_{k\in\partial i \setminus j}J_{ik}\mu_{k\setminus i}(t)C^{\rm dc}_{i\setminus j}(t,t') \nonumber\\
  & \quad + \sum_{k\in \partial i \setminus j}J_{ik}J_{ki}\int_0^t ds R_{k\setminus i}(t,s)\left(C^{\rm dc}_{i\setminus j}(t,s)\mu_{i\setminus j}(t') + \mu_{i\setminus j}(t)C^{\rm dc}_{i\setminus j}(s,t') + C^{\rm dc}_{i\setminus j}(t,t')\mu_{i\setminus j}(s)\right) \nonumber \\
  & \quad - 2\sum_{k\in \partial i \setminus j}J_{ik}J_{ki}\int_0^t ds R_{k\setminus i}(t,s)\mu_{i\setminus j}(t)\mu_{i\setminus j}(s)\mu_{i\setminus j}(t') + 2T \mathcal{L}_{i\setminus j}(t)R_{i\setminus j}(t',t) \nonumber \\
  & \quad + \sum_{k\in \partial i \setminus j} J_{ik}^2 \int ds C_{k\setminus i}(t,s) \left(\mu_{i\setminus j}(t) R_{i\setminus j}(t',s) + R_{i\setminus j}(t,s) \mu_{i\setminus j}(t') \right) + \lambda \mu_{i\setminus j}( t').\\
  \frac{\partial}{\partial t}R_{i\setminus j}(t,t') & = R_{i\setminus j}(t,t') \left(1 - 2\mu_{i\setminus j}(t)\right) + \sum_{k\in\partial i \setminus j}J_{ik}\mu_{k\setminus i}(t)R_{i\setminus j}(t,t') \nonumber\\
  & \quad + \sum_{k\in \partial i \setminus j} J_{ik}J_{ki} \int ds R_{k\setminus i}(t,s) \left(R_{i\setminus j}(t,t')\mu_{i\setminus j}(s)+\mu_{i\setminus j}(t) R_{i\setminus j}(s,t') \right).
\end{align}
where we have introduced the function
\begin{equation}
  \mathcal{L}_{i\setminus j}(t) = \begin{cases}
    0 & g(x) = 0 \quad \ \text{deterministic}\\
    1 & g(x) = x \quad \ \text{demographic}\\
    \mu_{i\setminus j}(t) & g(x) = x^2 \quad \text{environmental}.
  \end{cases}
\end{equation}

We can finally obtain the self-consistent equation for the connected cavity correlations by subtracting $\dot{\mu}_{i\setminus j}(t)\mu_{i\setminus j}(t)$ from the self-consistent equation for the disconnected cavity correlations:
\begin{align}
  \frac{\partial}{\partial t}C_{i\setminus j}(t,t') & = \left(1 - 2 \mu_{i\setminus j}(t) \right) C_{i\setminus j}(t,t') + \sum_{k\in\partial i \setminus j}J_{ik}\mu_{k\setminus i}(t)C_{i\setminus j}(t,t') \nonumber\\
  & \quad + \sum_{k\in \partial i \setminus j}J_{ik}J_{ki}\int_0^t ds R_{k\setminus i}(t,s)\left(\mu_{i\setminus j}(t)C_{i\setminus j}(s,t') + C_{i\setminus j}(t,t')\mu_{i\setminus j}(s)\right) \nonumber \\
  & \quad + 2T \mathcal{L}_{i\setminus j}(t)R_{i\setminus j}(t',t) + \sum_{k\in \partial i \setminus j} J_{ik}^2 \int_0^{{\max (t,t')}} ds C_{k\setminus i}(t,s) \left(\mu_{i\setminus j}(t) R_{i\setminus j}(t',s) + R_{i\setminus j}(t,s) \mu_{i\setminus j}(t') \right).
\end{align}


\section{Average over disorder vs small coupling expansion}
When the average over disordered interactions is taken into account, one has to deal with an integral of the form
\begin{equation}
  \overline{\int \mathcal{D}y \mathcal{D}\hat{y} \mathcal{D}w \mathcal{D}\hat{w}\,c[y,\hat{y},w,\hat{w}] \exp{\left[-{\rm i}\int dt \left(J\hat{w}(t)x(t) + J'\hat{f}(t) y(t) \right)\right]}},
\end{equation}
where $\overline{\cdots}$ denotes the average over disordered couplings $(J,J')$, i.e. $\overline{\cdots}=\int dJdJ'\,p(J,J')\cdots$, where $p(J,J')$ is a bivariate Gaussian probability distribution function with averages $\overline{J}=\overline{J'}=m/K$, variances  $\overline{(J-\overline{J})^2}=\overline{(J'-\overline{J'})^2}=\sigma^2/K$ and correlation $\overline{(J-\overline{J})(J-\overline{J})}=\gamma \sigma^2/K$.

Here we demonstrate that when dealing with a large connectivity expansion the average over disorder can be exchanged with the small-coupling expansion, provided one only retains terms of order $1/K$ in the expansion.

We use the following shorthand notation for the sake of shortness: $\mathcal{D}c=\mathcal{D}y \mathcal{D}\hat{y} \mathcal{D}w \mathcal{D}\hat{w}\,c[y,\hat{y},w,\hat{w}]$, $a\cdot b = \int dt\, a(t) b(t)$, with $a$ and $b$ being $x$, $y$, $\hat{f}$ and $\hat{w}$.

\subsection*{Average over disorder first}
Since the average over disorder involves a Gaussian integral, it can be easily computed as follows
\begin{align}
  \overline{\int \mathcal{D}c \, e^{-{\rm i}\left(J \hat{w}\cdot x+J'\hat{f}\cdot y\right)}}&=\int \mathcal{D}c \, \overline{e^{-{\rm i}\left(J \hat{w}\cdot x+J'\hat{f}\cdot y\right)}}\\
  &=\int \mathcal{D}c\,\exp\left[-\frac{m}{K}\left({\rm i}\hat{w}\cdot x+ {\rm i}\hat{f}\cdot y \right)-\frac{\sigma^2}{2K}\left( (\hat{w}\cdot x)^2+(\hat{f}\cdot y)^2 - 2\gamma ({\rm i}\hat{w}\cdot x) ({\rm i}\hat{f}\cdot y) \right) \right].
\end{align}
Expanding at first order in $1/K$, i.e. large connectivity expansion
\begin{align}
  \overline{\int \mathcal{D}c \, e^{-{\rm i}\left(J \hat{w}\cdot x+J'\hat{f}\cdot y\right)}}&=\int \mathcal{D}c \,\left[1 - \frac{m}{K}\left({\rm i}\hat{w}\cdot x+ {\rm i}\hat{f}\cdot y \right)-\frac{\sigma^2}{2K}\left( (\hat{w}\cdot x)^2+(\hat{f}\cdot y)^2 - 2\gamma ({\rm i}\hat{w}\cdot x) ({\rm i}\hat{f}\cdot y) \right)\right]\\
  &\approx 1 - \frac{m}{K}\int dt\, \left(\langle{\rm i} \hat{w}(t)\rangle x(t)+ {\rm i}\hat{f}(t) \langle y(t)\rangle \right)+\gamma\frac{\sigma^2}{K} \int dt dt'\,\langle y(t){\rm i}\hat{w}(t')\rangle {\rm i}\hat{f}(t)x(t')\nonumber\\
  &\quad - \frac{\sigma^2}{2K} \int dt dt'\,\left(\langle \hat{w}(t)\hat{w}(t')\rangle x(t)x(t')+\hat{f}(t)\hat{f}(t')\langle y(t)y(t')\rangle \right)\\
  &= 1 - \frac{m}{K}\int dt\, {\rm i}\hat{f}(t) \mu(t) +\gamma\frac{\sigma^2}{K} \int dt dt'\,{\rm i}\hat{f}(t) R(t,t') x(t') - \frac{\sigma^2}{2K} \int dt dt'\, \hat{f}(t)C^{\rm dc}(t,t')\hat{f}(t'),
\end{align}
where we have defined the average over the measure $c[y,\hat{y},w,\hat{x}]$ as $\langle \cdots \rangle$. We have introduced the average $\mu(t)=\langle y(t)\rangle$, the (disconnected) correlation function $C^{\rm dc}(t,t')=\langle y(t)y(t')\rangle$ and the response function $R(t,t')=\langle y(t){\rm i}\hat{w}(t')\rangle$. The averages of hatted quantities only vanish due to causality constraints.
By re-exponetiating we get the following approximate expression
\begin{align}
  \overline{\int \mathcal{D}c \, e^{-{\rm i}\left(J \hat{w}\cdot x+J'\hat{f}\cdot y\right)}} &\approx \exp\Bigg\{ - \frac{m}{K}\int dt\, {\rm i}\hat{f}(t) \mu(t) +\gamma\frac{\sigma^2}{K} \int dt dt'\,{\rm i}\hat{f}(t) R(t,t') x(t')\nonumber\\
  &\qquad - \frac{\sigma^2}{2K} \int dt dt'\, \hat{f}(t)C^{\rm dc}(t,t')\hat{f}(t') \Bigg\}.
\end{align}


\subsection*{Small coupling expansion first}
We expand the exponential at second order in $J$, i.e. small coupling expansion,
\begin{align}
  \overline{\int \mathcal{D}c \, e^{-{\rm i}\left(J \hat{w}\cdot x+J'\hat{f}\cdot y\right)}} & \approx\overline{\int \mathcal{D}c\,\left[1-\left(J {\rm i}\hat{w}\cdot x+J'{\rm i}\hat{f}\cdot y\right)+\frac{1}{2}\left(J {\rm i}\hat{w}\cdot x+J'{\rm i}\hat{f}\cdot y\right)^2 \right]}\\
  &=\overline{ 1 - \int dt\, \left(J \langle {\rm i}\hat{w}(t)\rangle x(t) + J' {\rm i} \hat{f}(t) \langle y(t)\rangle\right) + \frac{1}{2}\int dt dt' \, \Bigg(J^2 \langle {\rm i}\hat{w}(t){\rm i}\hat{w}(t')\rangle x(t)x(t')}\nonumber\\
  &\qquad \overline{+{J'}^2 {\rm i}\hat{f}(t){\rm i}\hat{f}(t')\langle y(t)y(t')\rangle +2JJ' \langle y(t) {\rm i}\hat{w}(t')\rangle {\rm i}\hat{f}(t)x(t')\Bigg)}\\
  &=\overline{ 1 - \int dt\, J' {\rm i} \hat{f}(t) \mu(t) + \frac{1}{2}\int dt dt' \, \Bigg({J'}^2 {\rm i}\hat{f}(t)C^{\rm dc}(t,t'){\rm i}\hat{f}(t') + 2JJ' {\rm i}\hat{f}(t) R(t,t') x(t')\Bigg)}.\label{eq:small-coupling_expanded}
\end{align}
By re-exponentiating, we have to carefully consider the square of first order terms, i.e. $1-\langle x\rangle+\langle x^2\rangle/2 \approx \exp(\langle x\rangle+\langle x^2\rangle/2 - \langle x \rangle^2/2 )$ when $x \approx 0$.
\begin{align}
  \overline{\int \mathcal{D}c \, e^{-{\rm i}\left(J \hat{w}\cdot x+J'\hat{f}\cdot y\right)}} & \approx \overline{ \exp \Bigg \{ - \int dt\, J' {\rm i} \hat{f}(t) \mu(t) + \frac{1}{2}\int dt dt' \, {J'}^2 {\rm i}\hat{f}(t)C(t,t'){\rm i}\hat{f}(t')}\nonumber\\
  &\qquad \overline{ +\int dt dt'\, JJ' {\rm i}\hat{f}(t) R(t,t') x(t')\Bigg\}},
\end{align}
where we have introduced the (connected) correlation function $C(t,t')=C^{\rm dc}(t,t')-\mu(t)\mu(t')$. We defined, for the sake of shortness, the terms $A=\int dt\, {\rm i}\hat{f}(t)\mu(t)$, $B=\int dt dt'\, {\rm i}\hat{f}(t)C(t,t'){\rm i}\hat{f}(t')$, and $C=\int dt dt'\, {\rm i}\hat{f}(t)R(t,t')x(t')$. By Gaussian integration, we obtain
\begin{align}
  \overline{\int \mathcal{D}c \, e^{-{\rm i}\left(J \hat{w}\cdot x+J'\hat{f}\cdot y\right)}} & \approx \overline{ \exp \left( -AJ'+\frac{B}{2}{J'}^2+CJJ'\right)}\\
  &=\frac{\exp \left[\frac{\sigma^2 \left(A^2 K^2+2 A C (\gamma-1) K m-2 C^2 (\gamma-1) m^2\right)+K m (m (B+2 C)-2 A K)}{2 K \left(-K \sigma^2 (B+2 C \gamma)+C^2 \left(\gamma^2-1\right) \sigma^4+K^2\right)}\right]}{\sqrt{1+\frac{C^2 \left(\gamma^2-1\right) \sigma^4-K \sigma^2 (B+2 C \gamma)}{K^2}}}.
\end{align}
Retaining only the first order terms in $1/K$, i.e. large connectivity expansion
\begin{align}
    \overline{\int \mathcal{D}c \, e^{-{\rm i}\left(J \hat{w}\cdot x+J'\hat{f}\cdot y\right)}} & \approx 1 - \frac{m}{K}A + \frac{\sigma^2}{2K}\left(A^2+B+2\gamma C\right)\\
    &\approx \exp\left[- \frac{m}{K}A + \frac{\sigma^2}{2K}\left(A^2+B+2\gamma C\right)\right]\\
    &=\exp \Bigg\{ -\frac{m}{K}\int dt\, {\rm i}\hat{f}(t)\mu(t)+\frac{\sigma^2}{2K}\int dtdt'\, {\rm i}\hat{f}(t)\left(\mu(t)\mu(t')+C(t,t')\right){\rm i}\hat{f}(t')\nonumber\\
    &\qquad + \gamma\frac{\sigma^2}{K}\int dt dt'\, {\rm i}\hat{f}(t)R(t,t')x(t')\Bigg\}\\
    &=\exp \Bigg\{ -\frac{m}{K}\int dt\, {\rm i}\hat{f}(t)\mu(t)+\frac{\sigma^2}{2K}\int dtdt'\, {\rm i}\hat{f}(t)C^{\rm dc}(t,t'){\rm i}\hat{f}(t')\nonumber\\
    &\qquad + \gamma\frac{\sigma^2}{K}\int dt dt'\, {\rm i}\hat{f}(t)R(t,t')x(t')\Bigg\},
\end{align}
which is the same approximate result obtained by taking first the average over disorder. We note that the same result is obtained if we take the average over disorder before re-exponentiating, i.e in Eq.~\eqref{eq:small-coupling_expanded}, 
\begin{align}
    \overline{\int \mathcal{D}c \, e^{-{\rm i}\left(J \hat{w}\cdot x+J'\hat{f}\cdot y\right)}} & \approx 1 - \int dt\, \overline{J'} {\rm i} \hat{f}(t) \mu(t) + \frac{1}{2}\int dt dt' \, \Bigg(\overline{{J'}^2} {\rm i}\hat{f}(t)C^{\rm dc}(t,t'){\rm i}\hat{f}(t') + 2\overline{JJ'} {\rm i}\hat{f}(t) R(t,t') x(t')\Bigg)\\
    &\approx 1 - \int dt\, \frac{m}{K} {\rm i} \hat{f}(t) \mu(t) + \frac{1}{2}\int dt dt' \, \Bigg(\frac{\sigma^2}{K} {\rm i}\hat{f}(t)C^{\rm dc}(t,t'){\rm i}\hat{f}(t') + 2\gamma\frac{\sigma^2}{K} {\rm i}\hat{f}(t) R(t,t') x(t')\Bigg)\\
    &\approx \exp \Bigg\{ -\frac{m}{K}\int dt\, {\rm i}\hat{f}(t)\mu(t)+\frac{\sigma^2}{2K}\int dtdt'\, {\rm i}\hat{f}(t)C^{\rm dc}(t,t'){\rm i}\hat{f}(t')\nonumber\\
    &\qquad + \gamma\frac{\sigma^2}{K}\int dt dt'\, {\rm i}\hat{f}(t)R(t,t')x(t')\Bigg\}.
\end{align}
\end{document}
